\subsection{EdgeColoring}
\begin{definition}\label{5.25}
    Let $k\in\mathbb{N}$. A (proper) \textbf{$k$-edge-coloring}\index{edge coloring} of a graph $G$ is a function $c:E(G)\to\{1,2,\ldots,k\}$ such that no two adjacent edges have the same color.\\
    The \textbf{chromatic index (edge-chromatic-number)}\index{chromatic index} $\chi'(G)$ of a graph $G$ is the smallest $k$ such that $G$ has a proper $k$-edge-coloring.
\end{definition}
\begin{remark}\label{5.26}
    $\Delta(G)\leq \chi'(G)$
\end{remark}
\begin{theorem}[Vizing's theorem]\label{5.27}
    Every graph $G$ has a proper edge coloring with at most $\Delta(G)+1$ colors. Thus $\chi'(G)\in\{\Delta(G),\Delta(G)+1\}$.
\end{theorem}
\begin{proof}
    not here.
\end{proof}